% Part: proof-theory
% Chapter: proof-search
% Section: tableaux

\documentclass[../../../include/open-logic-section]{subfiles}

\begin{document}

\olfileid{pt}{ps}{tab}

\olsection{Tableaux}

Tableaux systems are !!{proof} systems closely related to sequent
calculi that are particularly well-suited for proof search, both by
hand and implemented in software. Rather than build trees of sequents,
in tableaux !!a{proof} is a tree of !!{formula}s. But unlike natural
deduction, the rules mirror those of the sequent calculus. A tableau
!!{proof} is essentially an explicit search for a model. Hence they
are sometimes called \emph{semantic} tableaux or also \emph{truth
trees}.

A \emph{signed} tableau is a tree of signed !!{formula}s, which can be
of the form $\sFmla{\True}{!A}$ or $\sFmla{\False}{!A}$. Trees grow
downwards. They begin with !!{formula}s that specify what we want to
!!{prove}. E.g., to !!{prove} a sequent $!A, !B \Sequent !C, !D$ it
would begin with $\sFmla{\True}{!A}, \sFmla{\True}{!B},
\sFmla{\False}{!C}, \sFmla{\True}{!D}$. !!^a{structure} that makes
$!A$ and $!B$ true and $!C$ and $!D$ false is !!a{structure} that
shows that would show that $!A, !B \Sequent !C, !D$ is not valid. 

A leaf (bottommost) node of the tableau can be extended by branch
extension rules which add additional nodes to the same branch or add
two sibling nodes below, splitting the branch into two new branches.
The rules can be applied to any signed !!{formula} above it in the
same branch. The branch extension rules for the classical tableaux
system \Log{Tc} are given in \olref{tab:Tc}.

\subfile{rules-Tc}

As you can see, the rules are just those of \Log{G3c} upside-down, and
the signs $\True$ and $\False$ indicate whether !!a{formula} is
principal on the left or right side of a sequent, respectively.

A branch of a tableau is \emph{closed} if it contains both
$\sFmla{\True}{!A}$ and $\sFmla{\False}{!A}$. If every branch is
closed, we say the entire tableaux is closed. For instance, here is a
clause tableau for $!A \lor (!B \land !C) \Sequent (!A \lor !B) \land
(!A \lor !C)$, i.e., starting with $\sFmla{\True}{!A \lor (!B \land
!C)}$ and $\sFmla{\False}{(!A \lor !B) \land (!A \lor !C)}$:
\begin{oltableau}
  [\sFmla{\True}{\formula{A} \lor (\formula{B} \land \formula{C})},
    just=\TAss, checked
    [\sFmla{\False}{(\formula{A} \lor \formula{B}) \land
        (\formula{A} \lor \formula{C})}, just=\TAss, checked
      [\sFmla{\True}{\formula{A}}, just={\TRule{\True}{\lor}[1]}
        [\sFmla{\False}{\formula{A} \lor \formula{B}},
          just={\TRule{\False}{\land}[2]}, checked
          [\sFmla{\False}{\formula{A}},
            just={\TRule{\False}{\lor}[4]}
            [\sFmla{\False}{\formula{B}},
              just={\TRule{\False}{\lor}[4]}, close]
          ]
        ]
        [\sFmla{\False}{\formula{A} \lor \formula{C}},
          just={\TRule{\False}{\land}[2]}, checked
          [\sFmla{\False}{\formula{A}},
            just={\TRule{\False}{\lor}[4]}
            [\sFmla{\False}{\formula{C}},
              just={\TRule{\False}{\lor}[4]}, close]
          ]
        ]
      ]
      [\sFmla{\True}{\formula{B} \land \formula{C}},
        just={\TRule{\True}{\lor}[1]}, checked
        [\sFmla{\False}{\formula{A} \lor \formula{B}},
          just={\TRule{\False}{\land}[2]}, checked
          [\sFmla{\True}{\formula{B}},
            just={\TRule{\True}{\land}[3]}, move by=2
            [\sFmla{\True}{\formula{C}},
              just={\TRule{\True}{\land}[3]}
              [\sFmla{\False}{\formula{A}},
                just={\TRule{\False}{\lor}[4]}
                [\sFmla{\False}{\formula{B}},
                  just={\TRule{\False}{\lor}[4]},close
                ]
              ]
            ]
          ]
        ]
        [\sFmla{\False}{\formula{A} \lor \formula{C}},
          just={\TRule{\False}{\land}[2]}, checked
          [\sFmla{\True}{\formula{B}},
            just={\TRule{\True}{\land}[3]}, move by=2
              [\sFmla{\True}{\formula{C}},
                just={\TRule{\True}{\land}[3]}
                [\sFmla{\False}{\formula{A}},
                  just={\TRule{\False}{\lor}[4]}
                  [\sFmla{\False}{\formula{C}},
                  just={\TRule{\False}{\lor}[4]},close
                ]
              ]
            ]
          ]
        ]
      ]
    ]
  ]
\end{oltableau}

The checkmarks indicate that the corresponding rule has been applied
to !!a{formula} on all branches below it. The $\otimes$ symbol
indicates that a branch is closed. The line numbers after the rules
serve as indications which signed !!{formula}s the rule is applied to
(namely the !!{formula} in the same branch on the line indicated).

Proof search is easy in tableaux. We generate a tableau in stages. At
stage $0$, simply write down the starting !!{formula}s at the top. At
stage $n+1$, do the following for each leaf node: Find the first
(top-most) signed formula not yet checked off. Apply the branch
extension rule to that formula. Once the rule has been applied in
every branch that includes the signed formula in question, check it
off. (The example above is generated using this algorithm.)

Signed !!{formula}s of the form $\sFmla{\True}{\lforall}$ and
$\sFmla{\False}{\lexists}$ are never checked off, so must be treated
specially. If they are present, we must ensure that the
$\TRule{\True}{\forall}$ and $\TRule{\False}{\lexists}$ rules are
eventually applied for all applicable terms. This can be guaranteed,
for instance, by applying them at every stage to all such formulas on
every branch (after dealing with the topmost signed !!{formula} not of
this form). The term $t$ to use is the first term~$t$ that can be
built from the !!{constant}s already occurring on the branch (or $\Obj
c_0$ if no !!{constant} occurs) and the !!{function}s occurring in the
starting !!{formula}s, and such that the corresponding signed
!!{formula} $\sFmla{\True}{!B(t)}$ or~$\sFmla{\False}{!B(t)}$ does not
already occur on the branch.

If this proof search does not produce a closed tableau, it either
contains a finite open branch where all non-atomic !!{formula}s are
checked off, or the search produces an infinite, finitely branching
tree, which must contain an infinite branch. Just like in
\olref[cpl]{sec}, we can define !!a{structure}~$\Struct{M}$ such that
if $\sFmla{\True}{!A}$ occurs on the branch, then $\Sat{M}{!A}$ and if
$\sFmla{\False}{!A}$ occurs, then $\Sat/{M}{!A}$. (Take $\Theta =
\Setabs{!A}{\sFmla{\True}{!A} \text{ on the branch}}$ and $\Xi =
\Setabs{!A}{\sFmla{\False}{!A} \text{ on the branch}}$.)

Tableaux generated by this proof search method can easily be
translated into \Log{G3c} proofs by assigning a sequent to every node.
If the starting formulas correspond to a sequent $\Gamma \Sequent
\Delta$, label each of the starting formulas by $\Gamma \Sequent
\Delta$. If a branch is extended by applying a branch extension rule
to a signed !!{formula} $\sFmla{\True}{!A}$, then the leaf is
labelled $!A, \Pi \Sequent \Lambda$, and if it is applied to a signed
!!{formula} $\sFmla{\False}{!A}$, the leaf is labelled $\Pi \Sequent
\Lambda, !A$. Where the tableau applies a \TRule{\True}{\star} rule,
apply the \LeftR{\star} rule to the sequent (with $!A$ principal);
where it applies a \TRule{\False}{\star} rule, apply the \RightR{\star}
rule. The premises of the rule are the sequents that label the
corresponding nodes added to the tableau. Once a tableau closes with
$\sFmla{\True}{!A}$ and $\sFmla{\False}{!A}$, the label of the final
leaf node will be a sequent of the form $!A, \Pi \Sequent \Lambda,
!A$. Some successive nodes in the tableau will be assigned the same
sequent; delete duplicates. For instance, the tableau above translates
into the !!{proof}:
\[\small
\Axiom$!A \fCenter !A, !B$
\RightLabel{\RightR{\lor}}
\UnaryInf$!A \fCenter !A \lor !B$
\Axiom$!A \fCenter !A, !C$
\RightLabel{\RightR{\lor}}
\UnaryInf$!A \fCenter !A \lor !C$
\RightLabel{\RightR{\land}}
\BinaryInf$!A \fCenter (!A \lor !B) \land (!A \lor !C)$
\Axiom$!B, !C \fCenter !A, !B$
\RightLabel{\RightR{\lor}}
\UnaryInf$!B, !C \fCenter !A \lor !B$
\RightLabel{\LeftR{\land}}
\UnaryInf$!B \land !C \fCenter !A \lor !B$
\Axiom$!B, !C \fCenter !A, !C$
\RightLabel{\RightR{\lor}}
\UnaryInf$!B, !C \fCenter !A \lor !C$
\RightLabel{\LeftR{\land}}
\UnaryInf$!B \land !C \fCenter !A \lor !C$
\RightLabel{\RightR{\land}}
\BinaryInf$!B \land !C \fCenter (!A \lor !B) \land
(!A \lor !C)$
\RightLabel{\LeftR{\lor}}
\BinaryInf$!A \lor (!B \land !C) \fCenter (!A \lor !B) \land
(!A \lor !C)$
\DisplayProof
\]

\end{document}
